% ch04_u04 —— 原书 p.152–p.161（4.3.3–4.3.16）
% 承接 p.151 末 4.3.2 的证明（$f(g)=gL$）。
\SourcePage{152}

置 $\mathcal{P}':=\mathcal{P}\setminus\{\text{零截面}\}$；这是 $\operatorname{Lagr}(V)$ 上的 $\mathbb{G}_{m}$-torsor。
易证自然态射 $\tilde{O}(V)\to f^{*}\mathcal{P}'$ 是同构，故 $\tilde{O}(V)$ 是 $O(V)$ 上的 $\mathbb{G}_{m}$-torsor。

\emph{注.} 设 $L\in\operatorname{Lagr}(V)$。则 (193) 在 $O(V)$ 中 $L$ 的稳定化子上典范分裂：若 $g\in O(V)$、$gL=L$，则 $g$ 在 $\tilde{O}(V)$ 中有唯一原像，它在 $M^{L}$ 上作用为恒等。

\bdp{4.3.3} 置 $O:=\mathbb{C}[[t]]$，$K:=\mathbb{C}((t))$。以 $\omega_{O}$ 表 $O$ 的（完备化）微分模。固定 $\omega_{O}$ 的一个平方根，即一维自由 $O$-模 $\omega_{O}^{1/2}$，连同同构 $\omega_{O}^{1/2}\otimes\omega_{O}^{1/2}\xrightarrow{\ \sim\ }\omega_{O}$。$W$ 的含义同 4.3.1。我们要构造 $O_{n}(K):=O(W\otimes K)$ 的中心扩张，其中 $O_{n}(K)$ 看作 $\mathbb{C}$ 上的群归纳概形。

置 $\omega_{K}^{1/2}:=\omega_{O}^{1/2}\otimes_{O}K$，$\omega_{K}:=\omega_{O}\otimes_{O}K$。考虑 Tate 空间 $V:=\omega_{K}^{1/2}\otimes W$。$W$ 上的双线性型诱导 $K$-双线性型 $V\times V\to\omega_{K}$；与 $\operatorname{Res}:\omega_{K}\to\mathbb{C}$ 复合，得非退化对称双线性型 $V\times V\to\mathbb{C}$，其型为偶。把扩张 (193) 限制到 $O_{n}(K)\hookrightarrow O(V)$，得中心扩张
\begin{equation}\tag{194}
0\to\mathbb{G}_{m}\to\tilde{O}_{n}(K)\to O_{n}(K)\to 0.
\end{equation}
它在 $O_{n}(O)\subset O_{n}(K)$ 上典范分裂（用 4.3.2 末的注，取 $L=\omega_{O}^{1/2}\otimes W\subset V$）。群 $\operatorname{Aut}\omega_{O}^{1/2}=\mu_{2}$ 作用在扩张 (194) 上，并保持 $O_{n}(O)$ 上的分裂。

\bdp{4.3.4} \emph{引理.} 由 $-1\in\operatorname{Aut}\omega_{O}^{1/2}$ 诱导的 $\tilde{O}_{n}(K)$ 自同构，把 $\tilde{g}\in\tilde{O}_{n}(K)$ 映到 $(-1)^{\theta(g)}\tilde{g}$，其中 $g$ 是 $\tilde{g}$ 在 $O_{n}(K)$ 中的像，$\theta:O_{n}(K)\to K^{*}/(K^{*})^{2}=\mathbb{Z}/2\mathbb{Z}$ 是旋量范数。

\emph{证明.} 由 (192)，只需证 $g\in O_{n}(K)\subset O(V)$ 时 $\chi(g)=\theta(g)$。按 $\theta$ 的定义（见 [D71] 第 II 章 §7），只需证：若 $g$ 是关于某非迷向 $x\in K^{n}$ 的正交补的反射，则 $\chi(g)$ 等于 $(x,x)\in K^{*}$ 在 $K^{*}/(K^{*})^{2}=\mathbb{Z}/2\mathbb{Z}$ 中的像。可设 $x\in O^{n}$，$x\notin tO^{n}$。

\SourcePage{153}
$L:=\omega_{O}^{1/2}\otimes W$ 是 $V$ 中的 Lagrange c-格，故 $\chi(g)$ 是 $\dim L/(L\cap gL)=\dim O/(x,x)O$ 的奇偶性。

\emph{注记.}
\begin{enumerate}[label=(\roman*),leftmargin=2.2em]
\item 不用反射，也可计算 $\chi$ 在 $SO_{n}(K)$ 的某分裂嘉当子群上的限制，并注意 $g\in O_{n}(\mathbb{C})$ 时 $\chi(g)=0$。
\item $\theta$ 在 $SO_{n}(K)$ 上的限制是边缘态射
\begin{equation}\tag{195}
SO_{n}(K)\to H^{1}(K,\mu_{2})=\mathbb{Z}/2\mathbb{Z},
\end{equation}
它来自正合列 $0\to\mu_{2}\to\operatorname{Spin}_{n}\to SO_{n}\to 0$。
\item 若 $g\in O_{n}(K)=O(W\otimes K)$，则 $\dim(W\otimes O)/((W\otimes O)\cap g(W\otimes O))$ 为偶，当且仅当 $\theta(g)=0$。这由引理 4.3.4 的证明推出。
\end{enumerate}

\bdp{4.3.5} 考虑 (194) 在 $SO_{n}(K)$ 上的限制：
\begin{equation}\tag{196}
0\to\mathbb{G}_{m}\to\tilde{SO}_{n}(K)\to SO_{n}(K)\to 0.
\end{equation}
它在 $SO_{n}(O)$ 上典范分裂。扩张 (196) 依赖 $\omega_{O}^{1/2}$ 的选取，故更应写成 $\tilde{SO}_{n}(K)_{\mathcal{C}}$，其中 $\mathcal{C}$ 是 $\omega_{O}$ 的平方根。设 $\mathcal{C}'$ 是 $\omega_{O}$ 的另一平方根，则 $\mathcal{C}'=\mathcal{C}\otimes\mathcal{A}$，其中 $\mathcal{A}$ 是 $\operatorname{Spec} O$ 上的 $\mu_{2}$-torsor（或 $\operatorname{Spec}\mathbb{C}$ 上的，二者一样）。考虑 $\mathbb{Z}/2\mathbb{Z}$ 被 $\mathbb{G}_{m}$ 的（平凡）扩张，使 $\mathcal{A}$ 成为其分裂构成的 $\mu_{2}$-torsor。它经 (195) 的拉回是（平凡）扩张
\begin{equation}\tag{197}
0\to\mathbb{G}_{m}\to\tilde{SO}_{n}(K)_{\mathcal{A}}\to SO_{n}(K)\to 0,
\end{equation}
并带 $SO_{n}(O)$ 上的分裂（4.1.8 已在更一般的情形引入此扩张）。

引理 4.3.4 给出 $\tilde{SO}_{n}(K)_{\mathcal{C}'}$ 与扩张 $\tilde{SO}_{n}(K)_{\mathcal{C}}$、$\tilde{SO}_{n}(K)_{\mathcal{A}}$ 之和之间的典范同构。它与 $SO_{n}(O)$ 上的分裂相容。

\SourcePage{154}
\bdp{4.3.6} 设 $S$、$O_{S}$、$K_{S}$ 的含义同 4.1.7。固定 $\mathcal{L}\in\omega^{1/2}(X)$，记 $\omega_{K_{S}}^{1/2}$ 为 $\mathcal{L}$ 到 $\operatorname{Spec} K_{S}$ 的拉回的截面空间。然后照 4.3.3 的办法：置 $V:=\omega_{K_{S}}^{1/2}\otimes W$，用「留数之和」映射 $\omega_{K_{S}}\to\mathbb{C}$ 定义 $V$ 上的标量积，把 $SO_{n}(K_{S})$ 嵌入 $O(V)$，最终得中心扩张
\begin{equation}\tag{198}
0\to\mathbb{G}_{m}\to\tilde{SO}_{n}(K_{S})_{\mathcal{L}}\to SO_{n}(K_{S})\to 0,
\end{equation}
并带 $SO_{n}(O_{S})$ 上的典范分裂。

\emph{注.} (198) 是各扩张 (196)（$K=K_{x}$，$x\in S$）的「超和」。说明如下：设 $G_{i}$（$i\in I$）是带态射 $\theta_{i}:G_{i}\to\mathbb{Z}/2\mathbb{Z}$ 的群，$\tilde{G}_{i}$ 是 $G_{i}$ 被 $\mathbb{G}_{m}$ 的中心扩张，则这些扩张的超和是 $\bigoplus_{i\in I}G_{i}$ 被 $\mathbb{G}_{m}$ 的扩张：在通常的和上，添加上标准扩张
\[
0\to\mathbb{G}_{m}\to A\to\bigoplus_{i\in I}(\mathbb{Z}/2\mathbb{Z})\to 0
\]
的拉回。其中 $A$ 由 $\mathbb{G}_{m}$ 与元素 $e_{i}$（$i\in I$）生成，定义关系为 $e_{i}^{2}=1$、$ce_{i}=e_{i}c$（$c\in\mathbb{G}_{m}$）、$e_{i}e_{j}=(-1)\cdot e_{j}e_{i}$（$i\neq j$）。在我们的情形，$\theta_{x}:SO_{n}(K_{x})\to\mathbb{Z}/2\mathbb{Z}$ 是旋量范数。

若 $\mathcal{L},\mathcal{L}'\in\omega^{1/2}(X)$，则 $\mathcal{L}'=\mathcal{L}\otimes\mathcal{E}$，其中 $\mathcal{E}$ 是 $\mu_{2}$-torsor。由 4.3.5，$\tilde{SO}_{n}(K_{S})_{\mathcal{L}'}$ 与扩张 $\tilde{SO}_{n}(K_{S})_{\mathcal{L}}$、$\tilde{SO}_{n}(K_{S})_{\mathcal{E}}$ 之和之间有典范同构（$\tilde{SO}_{n}(K_{S})_{\mathcal{E}}$ 的定义见 4.1.8）。

\bdp{4.3.7} 4.3.1 在 $\operatorname{Bun}_{SO_{n}}$ 上定义了线丛 $\operatorname{Pf}_{\mathcal{L}}$（$\mathcal{L}\in\omega^{1/2}(X)$）。记 $\operatorname{Pf}_{\mathcal{L}}^{S}$ 为 $\operatorname{Pf}_{\mathcal{L}}$ 到 4.1.7 所定义概形 $\operatorname{Bun}_{SO_{n},S}$ 的拉回。有 $SO_{n}(O_{S})\times\mathbb{G}_{m}$ 在 $\operatorname{Pf}_{\mathcal{L}}^{S}$ 上的显然作用（$\lambda\in\mathbb{G}_{m}$ 乘以 $\lambda$）。我们要把它延拓为 $\tilde{SO}_{n}(K_{S})_{\mathcal{L}}$ 在 $\operatorname{Pf}_{\mathcal{L}}^{S}$ 上的作用，与 $SO_{n}(K_{S})$ 在 $\operatorname{Bun}_{SO_{n},S}$ 上的作用相容。

设 $u\in\operatorname{Bun}_{SO_{n},S}$，$\tilde{g}\in\tilde{SO}_{n}(K_{S})_{\mathcal{L}}$。记 $\mathcal{F}$、$\mathcal{F}'$ 为对应 $u$ 与 $gu$ 的 $SO(W)$-丛，其中 $g\in SO_{n}(K_{S})$ 是 $\tilde{g}$ 的像。须定义同构 $\operatorname{Pf}_{\mathcal{L},\mathcal{F}}\cong\operatorname{Pf}_{\mathcal{L},\mathcal{F}'}$，即同构

\SourcePage{155}
\[
\operatorname{Pf}(W_{\mathcal{F}}\otimes\mathcal{L})\xrightarrow{\ \sim\ }\operatorname{Pf}(W_{\mathcal{F}'}\otimes\mathcal{L}).
\]
由 4.2.13，只需构造同构 $\operatorname{Pf}(V;L_{+},L_{-})\xrightarrow{\ \sim\ }\operatorname{Pf}(V;L_{+},L'_{-})$，其中 $V$ 是 4.3.6 的 Tate 空间，$L_{+}=\omega_{O_{S}}^{1/2}\otimes W\subset V$，而 $L_{-},L'_{-}\subset V$ 是离散 Lagrange 子空间，满足 $L'_{-}=gL_{-}$。由 (166)，这等价于构造同构 $f:(M_{L_{-}})^{\otimes-1}\xrightarrow{\ \sim\ }(M_{gL_{-}})^{\otimes-1}$。定义 $f$ 为 $\tilde{g}\in\tilde{O}(V)$ 在 $M$ 上作用所诱导。

注意：$\lambda\in\mathbb{G}_{m}\subset\tilde{SO}_{n}(K_{S})_{\mathcal{L}}$ 在 $\operatorname{Pf}_{\mathcal{L}}^{S}$ 上乘以 $\lambda^{-1}$。

\bdp{4.3.8} 如 4.3.1 所述，目标是给函子 (190) 定义一个典范的 $\ell^{\operatorname{Spin}}$-仿射结构。也就是说，对 $\mathcal{L}\in\omega^{1/2}(X)$ 与 $X$ 上的 $\mu_{2}$-torsor $\mathcal{E}$，须定义同构
\begin{equation}\tag{199}
\operatorname{Pf}_{\mathcal{L}}\otimes\ell^{\operatorname{Spin}}_{\mathcal{E}}\xrightarrow{\ \sim\ }\operatorname{Pf}_{\mathcal{L}'},\qquad \mathcal{L}':=\mathcal{L}\otimes\mathcal{E}.
\end{equation}
还须验证同构 (199) 的某些相容性质。

为简化记号，把 $\ell^{\operatorname{Spin}}_{\mathcal{E}}$ 写成 $\ell_{\mathcal{E}}$。设 $S\subset X$ 有限。4.1.7–4.1.8 构造了中心扩张 $\tilde{SO}_{n}(K_{S})_{\mathcal{E}}$ 在 $\ell^{S}_{\mathcal{E}}:=\ell_{\mathcal{E}}$ 到 $\operatorname{Bun}_{SO_{n},S}$ 的拉回上的作用。故由 4.3.6–4.3.7，$\tilde{SO}_{n}(K_{S})_{\mathcal{L}'}$ 同时作用在 $\operatorname{Pf}_{\mathcal{L}}^{S}\otimes\ell^{S}_{\mathcal{E}}$ 与 $\operatorname{Pf}_{\mathcal{L}'}^{S}$ 上。回忆 (199) 两端在平凡 $SO_{n}$-丛上的纤维都等于 $\mathbb{C}$。

\bdp{4.3.9} \emph{定理.} 存在唯一的同构 (199)，使对每个 $S$，相应同构 $\operatorname{Pf}_{\mathcal{L}}^{S}\otimes\ell^{S}_{\mathcal{E}}\xrightarrow{\ \sim\ }\operatorname{Pf}_{\mathcal{L}'}^{S}$ 是 $\tilde{SO}_{n}(K_{S})_{\mathcal{L}'}$-等变的，且由 (199) 诱导的平凡 $SO_{n}$-丛上纤维之间的同构是恒等。

证明见 4.3.11–4.3.13。较弱陈述的简短证明见 [BLaSo] 的 §5.2。

\bdp{4.3.10} \emph{命题.} 同构 (199) 给出函子 $\operatorname{Pf}:\omega^{1/2}(X)\to\nu^{1/2}(\operatorname{Bun}_{SO_{n}})$ 上的 $\ell^{\operatorname{Spin}}$-仿射结构。

证明见 4.3.15。

\SourcePage{156}
\bdp{4.3.11} 开始证定理 4.3.9。$n>2$ 时 (199) 的唯一性显然：此时 $SO_{n}$ 半单，故有 $G=SO_{n}$、$S\neq\emptyset$ 的同构 (155)。若 $n=2$，$SO_{n}(K_{S})$ 在 $\operatorname{Bun}_{SO_{n},S}$ 上的作用不可迁；但 adele 上的 $SO_{n}$ 在 $\varprojlim_{S}\operatorname{Bun}_{SO_{n},S}(\mathbb{C})$ 上可迁，这足以给出唯一性。

证存在性时设 $n>2$。$n=2$ 的情形可用嵌入 $SO_{2}\hookrightarrow SO_{3}$ 与相应态射 $\operatorname{Bun}_{SO_{2}}\to\operatorname{Bun}_{SO_{3}}$ 处理，或用 4.3.14 末的注处理。

考虑 $\operatorname{Bun}_{SO_{n},S}$ 上 $SO_{n}(K_{S})$-等变的线丛 $\mathcal{C}_{S}:=\operatorname{Pf}_{\mathcal{L}}^{S}\otimes\ell^{S}_{\mathcal{E}}\otimes(\operatorname{Pf}_{\mathcal{L}'}^{S})^{*}$。$\operatorname{Bun}_{SO_{n},S}$ 中对应平凡 $SO_{n}$-丛（连同 $S$ 上的显然平凡化）的点的稳定化子等于 $SO_{n}(A_{S})$，其中 $A_{S}:=H^{0}(X\setminus S,\mathcal{O}_{X})$。于是 $SO_{n}(K_{S})$ 在 $\mathcal{C}_{S}$ 上的作用诱导态射 $f_{S}:SO_{n}(A_{S})\to\mathbb{G}_{m}$。只需证所有 $S$ 的 $f_{S}$ 都平凡（则 $S\neq\emptyset$ 时可用 (155) 得 $\mathcal{C}_{S}$ 的 $SO_{n}(K_{S})$-等变平凡化；这些平凡化当然彼此相容）。

记 $\Sigma$ 为 $X$ 的有限子概形构成的概形（即 $X$ 各对称幂的不交并）。对不必约化的\footnote{这一点在 $S$ 变动时很重要。对固定的 $S$，环 $A_{S}$、$O_{S}$、$K_{S}$ 只依赖 $S_{\mathrm{red}}$。} $S\in\Sigma$，$A_{S}$、$O_{S}$、$K_{S}$ 有意义（例如 $O_{S}$ 是 $X$ 沿 $S$ 的完备化上的函数环），且环 $A_{S}$、$O_{S}$、$K_{S}$ 自然地组织成族（即显然可定义 $\Sigma$ 上三个环归纳概形，它们在 $S\in\Sigma$ 上的纤维分别为 $A_{S}$、$O_{S}$、$K_{S}$）。

易证这些态射 $f_{S}$ 构成族（即它们来自 $\Sigma$ 上群归纳概形之间的态射）。显然 $S\subset S'$ 时 $f_{S'}$ 到 $SO_{n}(A_{S})$ 的限制等于 $f_{S}$。4.3.12–4.3.13 将由此两点推出 $f_{S}=1$。

\SourcePage{157}
\bdp{4.3.12} 设 $Y$ 是 $\mathbb{C}$ 上有限型的分离概形，$R$ 是 $\mathbb{C}$-代数。置
\[
Y_{rat}(R)=\varinjlim_{U}\operatorname{Mor}(U,Y),
\]
其中极限取遍所有开 $U\subset X\otimes R$，使 $U$ 在 $\operatorname{Spec} R$ 的任一点上的纤维非空。换言之，$Y_{rat}(R)$ 的元素是由 $\operatorname{Spec} R$ 参数化的有理映射族 $X\to Y$。函子 $Y_{rat}$ 称为有理映射空间 $X\to Y$。易证 $Y_{rat}$ 是 fppf 拓扑的层，即 [LMB93] 意义下的「空间」。

有空间 $Y(A_{S})$（$S\in\Sigma$），它们构成族（即 $\Sigma$ 上有自然空间，其每个 $S$ 上的纤维等于 $Y(A_{S})$）。故 $Y_{rat}$ 上的正则函数定义 $Y(A_{S})$ 上的正则函数族 $f_{S}$（$S\in\Sigma$），使 $S\subset S'$ 时 $f_{S'}$ 到 $Y(A_{S})$ 的拉回等于 $f_{S}$。易见 $Y_{rat}$ 上的函数等同于具有此性质的函数族 $f_{S}$。

\bdp{4.3.13} \emph{命题.} 设 $G$ 是连通代数群。
\begin{enumerate}[label=(\roman*),leftmargin=2.2em]
\item $G_{rat}$ 上每个正则函数是常数。特别地，每个群态射 $G_{rat}\to\mathbb{G}_{m}$ 平凡。
\item 此外，对每个 $\mathbb{C}$-代数 $R$，$G_{rat}\otimes R$ 上每个正则函数是常数（即 $R$ 的元素）。
\end{enumerate}

\emph{证明.} 把 $G$ 表为 $\bigcup_{i}U_{i}$，其中 $U_{i}$ 是开集，同构于 $(\mathbb{A}^{1}\setminus\{0\})^{r}\times\mathbb{A}^{s}$（例如取 $U\subset G$ 为关于某 Borel 子群的大胞腔，则 $G$ 被有限多个形如 $gU$（$g\in G$）的集合覆盖）。有开覆盖 $G_{rat}=\bigcup_{i}(U_{i})_{rat}$，且 $(U_{i})_{rat}\cap(U_{j})_{rat}\neq\emptyset$。故只需对 $G=(\mathbb{G}_{m})^{r}\times(\mathbb{G}_{a})^{s}$ 证明命题。此外，只需对 $\mathbb{G}_{a}$ 与 $\mathbb{G}_{m}$ 证明 (ii)。

例如考虑 $\mathbb{G}_{m}$ 的情形。取 $X$ 上的丰沛线丛 $\mathcal{A}$，置 $V_{n}:=H^{0}(X,\mathcal{A}^{\otimes n})$，$V'_{n}:=V_{n}\setminus\{0\}$。定义 $\pi_{n}:V'_{n}\times V'_{n}\to(\mathbb{G}_{m})_{rat}$ 为 $(f,g)\mapsto f/g$。$(\mathbb{G}_{m})_{rat}\otimes R$ 上的正则函数 $\phi$ 定义 $(V'_{n}\times V'_{n})\otimes R$ 上的正则函数 $\pi_{n}^{*}\phi$，它关于 $\mathbb{G}_{m}$ 在 $V'_{n}\times V'_{n}$ 上的显然作用不变。$n$ 充分大时 $\dim V_{n}>1$，故 $\pi_{n}^{*}\phi$ 延拓为

\SourcePage{158}
$(V_{n}\times V_{n})\otimes R$ 上 $\mathbb{G}_{m}$-不变的正则函数，它必为常数。故 $\phi$ 为常数。

\bdp{4.3.14} 本小节后文不用（除 5.3.10 所需的 $\operatorname{GRAS}_{G}$ 的定义外）。

设 $G$ 是连通代数群。下述接近 $\operatorname{Bun}_{G}$ 的方式似乎自然。

记 $\operatorname{GRAS}_{G}$ 为 $X$ 上带有理截面的 $G$-torsor 构成的空间。其精确定义与 4.3.12 中 $Y_{rat}$ 的定义颇相似。我们愿称 $\operatorname{GRAS}_{G}$ 为对应 $G$ 与 $X$ 的\emph{大 Grassmann 流形}：因为对固定的有限 $S\subset X$，在 $X\setminus S$ 上平凡化的 $X$ 上 $G$-丛的空间可等同于归纳概形 $G(K_{S})/G(O_{S})=\prod_{x\in X}G(K_{x})/G(O_{x})$（见 5.3.10），而 $G(K_{x})/G(O_{x})$ 称为仿射 Grassmann 流形或环路 Grassmann 流形（见 4.5 或 [MV]）。

态射 $\pi:\operatorname{GRAS}_{G}\to\operatorname{Bun}_{G}$ 是光滑拓扑下的 $G_{rat}$-torsor（若 $G$ 的约化部分等于 $GL_{n}$、$SL_{n}$ 或 $Sp_{n}$，则存在光滑满射 $S\to\operatorname{Bun}_{G}$ 及截面 $S\to\operatorname{GRAS}_{G}$，这是显然的；一般 $G$ 可用 [DSim]）。

考虑函子
\begin{equation}\tag{200}
\pi^{*}:\operatorname{Vect}(\operatorname{Bun}_{G})\to\operatorname{Vect}(\operatorname{GRAS}_{G}),
\end{equation}
其中 $\operatorname{Vect}$ 表向量丛范畴。由 4.3.13，(200) 全忠实。可证对任意概形 $T$，$G_{rat}\times T$ 上每个向量丛都来自 $T$。这说明 (200) 是等价。

\emph{注.} (199) 的构造可解释如下：先构造 (199) 左右两端到 $\operatorname{GRAS}_{SO_{n}}$ 的拉回之间的同构，再用 (200) 的全忠实性。其实不必用同构 (155)。故 (199) 的构造在 $SO_{2}$ 的情形也有效。

\SourcePage{159}
\bdp{4.3.15} 证明命题 4.3.10。各同构 (199) 彼此相容（用 4.3.9 的唯一性）。剩下要证 (199) 的张量平方等于复合
\begin{equation}\tag{201}
\operatorname{Pf}_{\mathcal{L}}^{\otimes 2}\xrightarrow{\ \sim\ }\nu_{\mathcal{L}}\xrightarrow{\ \sim\ }\nu\xrightarrow{\ \sim\ }\nu_{\mathcal{L}'}\xrightarrow{\ \sim\ }\operatorname{Pf}_{\mathcal{L}'}^{\otimes 2},
\end{equation}
其中 $\nu_{\mathcal{L}}$ 由 (188) 定义。

固定 $X$ 上的 $SO_{n}$-torsor $\mathcal{F}$ 及其在某非空有限 $S\subset X$ 上的 $X\setminus S$ 平凡化。用该平凡化计算 (199) 与 (201) 诱导的同构 $\operatorname{Pf}_{\mathcal{L},\mathcal{F}}^{\otimes 2}\xrightarrow{\ \sim\ }\operatorname{Pf}_{\mathcal{L}',\mathcal{F}}^{\otimes 2}$。

回忆 $\operatorname{Pf}_{\mathcal{L},\mathcal{F}}:=\operatorname{Pf}(W_{\mathcal{F}}\otimes\mathcal{L})\otimes\operatorname{Pf}(W\otimes\mathcal{L})^{\otimes-1}$。由 4.2.13，
\[
\operatorname{Pf}(W_{\mathcal{F}}\otimes\mathcal{L})=\operatorname{Pf}(V;L_{+},L_{-}),\qquad
\operatorname{Pf}(W\otimes\mathcal{L})=\operatorname{Pf}(V;L_{+}^{0},L_{-}),
\]
其中 $V=\mathcal{L}_{K_{S}}\otimes W$，$L_{-}=\Gamma(X\setminus S,\mathcal{L}\otimes W)$，$L_{+}^{0}=\mathcal{L}_{O_{S}}\otimes W$，而 $L_{+}$ 是 $W_{\mathcal{F}}\otimes\mathcal{L}$ 到 $\operatorname{Spec} O_{S}$ 的拉回的截面空间（用 4.3.6 的记号）。用 (166) 得
\begin{equation}\tag{202}
\operatorname{Pf}_{\mathcal{L},\mathcal{F}}=M^{L_{+}}\otimes(M^{L_{+}^{0}})^{*},
\end{equation}
其中 $M$ 是 $\operatorname{Cl}(V)$ 上不可约的 $\mathbb{Z}/2\mathbb{Z}$-分次离散模。$\operatorname{Pf}_{\mathcal{L}',\mathcal{F}}$ 有类似的描述，用 $V'$、$L'_{+}$、$(L_{+}^{0})'$、$L'_{-}$，其中 $V'=\mathcal{L}'_{K_{S}}\otimes W$，等等。固定 4.3.8 的 $\mu_{2}$-torsor $\mathcal{E}$ 在 $S$ 上的平凡化。它给出 $\mathcal{E}$ 在 $\operatorname{Spec} O_{S}$ 上的平凡化，从而给出等同
\begin{equation}\tag{203}
(V',L'_{+},(L_{+}^{0})')=(V,L_{+},L_{+}^{0}).
\end{equation}
因 (202) 不含 $L_{-}$，得同构 $\operatorname{Pf}_{\mathcal{L},\mathcal{F}}\xrightarrow{\ \sim\ }\operatorname{Pf}_{\mathcal{L}',\mathcal{F}}$。易证它与 (199) 诱导者一致（注意 $\mathcal{F}$ 在 $X\setminus S$ 上的平凡化与 $\mathcal{E}$ 在 $S$ 上的平凡化给出 $\ell^{\operatorname{Spin}}_{\mathcal{E}}$ 在 $\mathcal{F}$ 上的平凡化，因为 (150) 左端有特异元素）。

现在须证 (201) 诱导的同构 $\operatorname{Pf}_{\mathcal{L},\mathcal{F}}^{\otimes 2}\xrightarrow{\ \sim\ }\operatorname{Pf}_{\mathcal{L}',\mathcal{F}}^{\otimes 2}$，在把 $\operatorname{Pf}_{\mathcal{L},\mathcal{F}}$、$\operatorname{Pf}_{\mathcal{L}',\mathcal{F}}$ 与 (202) 右端等同后是恒等。

\SourcePage{160}
$\mathcal{F}$ 在 $X\setminus S$ 上的平凡化给出同构 $\nu_{\mathcal{L}}(\mathcal{F})\xrightarrow{\ \sim\ }d(L_{+}^{0},L_{+})$，其中 $d(L_{+}^{0},L_{+})$ 是相对行列式，即对任意 c-格 $U\subset L\cap L_{+}^{0}$ 有 $d(L_{+}^{0},L_{+})=\det(L_{+}/U)\otimes\det(L_{+}^{0}/U)^{\otimes-1}$。类似有 $\nu_{\mathcal{L}'}(\mathcal{F})=d((L_{+}^{0})',L'_{+})$。(201) 中的同构 $\nu_{\mathcal{L}}(\mathcal{F})\xrightarrow{\ \sim\ }\nu_{\mathcal{L}'}(\mathcal{F})$ 在 [Del87] 中如下定义：任取 $\mathcal{L}$ 与 $\mathcal{L}'$ 到 $\operatorname{Spec} O_{S}$ 的拉回之间的同构 $f$。$f$ 给出同构 $f_{*}:(V,L_{+},L_{+}^{0})\xrightarrow{\ \sim\ }(V',L'_{+},(L_{+}^{0})')$，从而给出 $d(L_{+}^{0},L_{+})\xrightarrow{\ \sim\ }d(L'_{+},(L_{+}^{0})')$，它其实不依赖 $f$ 的选取。用 $\mu_{2}$-torsor $\mathcal{E}=\mathcal{L}'\otimes\mathcal{L}^{\otimes-1}$ 在 $\operatorname{Spec} O_{S}$ 上的上述平凡化来定义 $f$ 较方便。此时 $f_{*}$ 与 (203) 一致。

于是我们把 $\nu_{\mathcal{L}}(\mathcal{F})$ 与 $\nu_{\mathcal{L}'}(\mathcal{F})$ 都等同于 $d(L_{+}^{0},L_{+})$，使 (201) 中的同构成为恒等。又把 $\operatorname{Pf}_{\mathcal{L},\mathcal{F}}$ 与 $\operatorname{Pf}_{\mathcal{L}',\mathcal{F}}$ 都等同于 (202) 右端。剩下要证 (187) 的同构与 $\mathcal{L}'$ 的类似物诱导同一同构
\begin{equation}\tag{204}
(M^{L_{+}}\otimes(M^{L_{+}^{0}})^{*})^{\otimes 2}\xrightarrow{\ \sim\ }d(L_{+}^{0},L_{+}).
\end{equation}

据 4.2.8 与 4.2.13，(187) 诱导的 (204) 可描述如下。有典范同构
\begin{equation}\tag{205}
N^{L_{+}\otimes H}\otimes(N^{L_{+}^{0}\otimes H})^{*}\xrightarrow{\ \sim\ }d(L_{+}^{0},L),
\end{equation}
其中 $N$ 是 Clifford 代数 $\operatorname{Cl}(V\oplus V^{*})=\operatorname{Cl}(V\oplus V)=\operatorname{Cl}(V\otimes H)$ 上不可约的 $(\mathbb{Z}/2\mathbb{Z})$-分次离散模，$H$ 表 $\mathbb{C}^{2}$ 配双线性型 $\bigl(\begin{smallmatrix}0&1\\1&0\end{smallmatrix}\bigr)$（构造 (205) 时取 $N$ 为 (182) 的右端）。
% [Tutor Note: 原书 p.160 (205) 右端作 d(L_+^0,L)，而下文的 (204) 作 d(L_+^0,L_+)；此处照录原书。]
另一方面，$P:=M\otimes M$ 是 $\operatorname{Cl}(V)\otimes\operatorname{Cl}(V)=\operatorname{Cl}(V\otimes H'')$ 上不可约的 $(\mathbb{Z}/2\mathbb{Z})$-分次离散模，其中 $H''$ 表 $\mathbb{C}^{2}$ 配双线性型 $\bigl(\begin{smallmatrix}1&0\\0&1\end{smallmatrix}\bigr)$。把 (204) 左端改写为 $P^{L_{+}\otimes H''}\otimes(P^{L_{+}^{0}\otimes H''})^{*}$。于是正交同构 $\psi:H''\xrightarrow{\ \sim\ }H$ 诱导 (204)。要使 (204) 是 (187) 诱导者，须用 $\det\psi=i$ 规范化 $\psi$（还是 $-i$？？我们应当查一下！）。
% [Tutor Note: 原书 p.160 此句末尾的括注「还是 -i（原文用减号）？？我们应当查一下！」是作者在底本中的自注，照录。]

\SourcePage{161}
因上述描述不含 $L_{-}$，$\mathcal{L}'$ 的 (187) 类似物诱导同一同构 (204)，证毕。

\bdp{4.3.16} 本小节与 4.3.17 将用于 4.4.14（水平性定理 2.7.3 证明的末段）与定理 5.4.5（本工作的主要结果）的证明。读者暂时可以跳过。

照例置 $O:=\mathbb{C}[[t]]$，$K:=\mathbb{C}((t))$。固定 $\mathcal{L}\in\omega^{1/2}(X)$，即 $\mathcal{L}$ 是 $\omega_{X}$ 的平方根。再固定 $\omega_{O}$ 的一个平方根，记为 $\omega_{O}^{1/2}$。于是可定义 2.6.5 中概形 $X^{\wedge}$ 的二重覆叠 $X_{2}^{\wedge}$。回忆 $X^{\wedge}$ 的 $R$-点是一个 $R$-态射 $\alpha:\operatorname{Spec}(R\hat{\otimes}O)\to X\otimes R$，其微分在 $\operatorname{Spec} R$ 的任一点都不为零。记 $\mathcal{L}_{R}$ 为 $\mathcal{L}$ 到 $X\otimes R$ 的拉回。按定义，$X_{2}^{\wedge}(R)$ 在 $\gamma\in X^{\wedge}(R)$ 上的纤维，是 $R\hat{\otimes}\omega_{O}$ 的平方根群胚中同构 $H^{0}(\operatorname{Spec} R\hat{\otimes}O,\alpha^{*}\mathcal{L}_{R})\xrightarrow{\ \sim\ }R\hat{\otimes}\omega_{O}^{1/2})$ 的集合。
% [Tutor Note: 原书 p.161 此行末尾 $\omega_O^{1/2}$ 之后多一个右括号，括号不配对，属排版笔误，照录。]

3.5.2 引入的群归纳概形 $\operatorname{Aut}_{2}O:=\operatorname{Aut}(O,\omega_{O}^{1/2})$ 以结构搬运的方式作用在 $X_{2}^{\wedge}$ 上。

设 $M$ 是 2.8.1 中 $G=SO(W)=SO_{n}$ 这一特殊情形的概形。记 $M_{2}^{\wedge}$ 为 $M$ 与 $X_{2}^{\wedge}$ 在 $X$ 上的纤维积（故 $M_{2}^{\wedge}$ 是 2.8.3 中概形 $M^{\wedge}$ 的二重覆叠）。则半直积 $\operatorname{Aut}_{2}O\ltimes SO_{n}(K)$ 作用在 $M_{2}^{\wedge}$ 上。事实上 $M_{2}^{\wedge}$ 是 $M^{\wedge}$ 与 $X_{2}^{\wedge}$ 在 $X^{\wedge}$ 上的纤维积，而 $\operatorname{Aut}_{2}O\ltimes SO_{n}(K)$ 作用在图表
\[
\begin{tikzcd}
M^{\wedge} \arrow[d]\\
X_{2}^{\wedge} \arrow[r] & X^{\wedge}
\end{tikzcd}
\]
上（$\operatorname{Aut}O\ltimes SO_{n}(K)$ 在 $M^{\wedge}$ 上的作用在 2.8.4 中定义；$\operatorname{Aut}_{2}O\ltimes SO_{n}(K)$ 经其商 $\operatorname{Aut}_{2}O$ 与 $\operatorname{Aut}O$ 作用在 $X_{2}^{\wedge}$ 与 $X^{\wedge}$ 上）。

记 $\operatorname{Pf}_{\mathcal{L}}^{\wedge}$ 为 4.3.1 定义的 $\operatorname{Bun}_{SO_{n}}$ 上线丛 $\operatorname{Pf}_{\mathcal{L}}$ 到 $M_{2}^{\wedge}$ 的拉回。要把 $\operatorname{Aut}_{2}O\ltimes SO_{n}(K)$ 在 $M_{2}^{\wedge}$ 上的作用提升为 $\operatorname{Aut}_{2}O\ltimes\tilde{SO}_{n}(K)$ 在 $\operatorname{Pf}_{\mathcal{L}}^{\wedge}$ 上的作用，其中 $\tilde{SO}_{n}(K)$ 是中心扩张
% 原书 p.161 末句未完，续 p.162（下一单元）。
